\documentclass[a4paper,12pt]{ctexart}
\usepackage{xcolor,graphicx,amsmath}
\usepackage[top=1.8cm, bottom=1.8cm, left=1.8cm, right=1.8cm]{geometry}
\usepackage{array,hyperref}
\hypersetup{colorlinks=true,linkcolor=blue!70!black}
\linespread{1.35}

\definecolor{defblue}{RGB}{0,80,160}\definecolor{thinkorange}{RGB}{230,120,20}
\definecolor{formulared}{RGB}{180,30,50}\definecolor{funboxbg}{RGB}{245,248,252}
\definecolor{examplebg}{RGB}{255,250,240}\definecolor{practicegreen}{RGB}{40,130,70}

\newenvironment{thinkbox}{\vspace{0.3cm}\begin{center}\begin{minipage}{0.88\textwidth}\colorbox{funboxbg}{\begin{minipage}{0.94\textwidth}\vspace{0.15cm}\textbf{\textcolor{thinkorange}{【 想一想 】}}}{\vspace{0.15cm}\end{minipage}}\end{minipage}\end{center}\vspace{0.3cm}}
\newenvironment{definition}[1]{\vspace{0.2cm}\begin{center}\begin{minipage}{0.88\textwidth}\fbox{\begin{minipage}{0.92\textwidth}\vspace{0.1cm}\textbf{\textcolor{defblue}{【 #1 】}}}{\vspace{0.1cm}\end{minipage}}\end{minipage}\end{center}\vspace{0.2cm}}
\newenvironment{example}{\vspace{0.2cm}\begin{center}\begin{minipage}{0.88\textwidth}\colorbox{examplebg}{\begin{minipage}{0.94\textwidth}\vspace{0.15cm}\textbf{\textcolor{orange!80!black}{【 例题 】}}}{\vspace{0.15cm}\end{minipage}}\end{minipage}\end{center}\vspace{0.2cm}}
\newenvironment{practice}{\vspace{0.3cm}\begin{center}\begin{minipage}{0.88\textwidth}\colorbox{funboxbg}{\begin{minipage}{0.94\textwidth}\vspace{0.15cm}\textbf{\textcolor{practicegreen}{【 练一练 】}}}{\vspace{0.15cm}\end{minipage}}\end{minipage}\end{center}\vspace{0.2cm}}
\newenvironment{figplaceholder}[1]{\vspace{0.2cm}\begin{center}\fbox{\begin{minipage}{0.82\textwidth}\centering\textcolor{blue!70!black}{\textbf{【插图位置】}}\\[0.2cm]#1\end{minipage}}\end{center}\vspace{0.2cm}}
\newcommand{\formula}[1]{\begin{center}\textcolor{formulared}{\large\boxed{#1}}\end{center}}

\title{\textcolor{blue!70!black}{\Huge\textbf{化学2 · 常见元素与无机反应}}}
\author{}\date{}

\begin{document}
\maketitle\thispagestyle{empty}
\section*{\textcolor{blue!70!black}{致同学}}
化学1教了你化学的"语言"——符号、化学式、方程式、摩尔。现在我们要用这门语言来认识\textbf{身边的具体物质}：身边的空气、喝的水、用的金属、吃的盐。从这一本开始，化学不再是纸上的符号——而是你每天接触的实物。

\newpage
\section{\textcolor{orange!80!black}{第一课\quad 空气与氧气}}

\subsection*{空气的组成}

空气不是一种单一物质——它是混合物。干燥空气按体积分数：\textbf{氮气N$_2$约78\%，氧气O$_2$约21\%，稀有气体约0.94\%，二氧化碳CO$_2$约0.03\%}。

\subsection*{氧气——生命的"助燃剂"}

氧气的化学性质非常活泼——它能和大多数物质发生氧化反应。

\textbf{氧气的检验：}用带火星的木条伸入集气瓶中——木条\textbf{复燃}（重新燃烧起来），证明是氧气。

\textbf{实验室制取氧气：}三种常见方法：
\begin{enumerate}
  \item 加热高锰酸钾：2KMnO$_4 \xrightarrow{\Delta}$ K$_2$MnO$_4$ + MnO$_2$ + O$_2$$\uparrow$
  \item 过氧化氢分解（二氧化锰催化）：2H$_2$O$_2 \xrightarrow{\text{MnO}_2}$ 2H$_2$O + O$_2$$\uparrow$
  \item 加热氯酸钾（二氧化锰催化）：2KClO$_3 \xrightarrow[\text{MnO}_2]{\Delta}$ 2KCl + 3O$_2$$\uparrow$
\end{enumerate}

\subsection*{燃烧——剧烈的氧化反应}

\begin{definition}{燃烧}
可燃物与氧气发生的一种发光、放热的剧烈的\textbf{氧化反应}。
\end{definition}

\textbf{燃烧的三个条件（缺一不可）：}可燃物、氧气（或空气）、达到着火点（可燃物燃烧所需的最低温度）。

\textbf{灭火的原理（破坏任意一个条件）：}清除可燃物、隔绝氧气、降温至着火点以下。

\begin{example}
铁丝在氧气中剧烈燃烧的化学方程式是什么？现象是什么？

\textbf{解：}3Fe + 2O$_2 \xrightarrow{\text{点燃}}$ Fe$_3$O$_4$。现象：\textbf{火星四射}、生成\textbf{黑色固体}、放出大量热。实验时集气瓶底部要放少量水——防止高温熔化物溅落炸裂瓶底。
\end{example}

\begin{thinkbox}
为什么煤炉在冬天特别危险？煤炭不完全燃烧产生\textbf{一氧化碳CO}——无色无味，和血红蛋白的结合力比氧气强200多倍。人吸入CO后，血液"装满了CO而没装氧气"——导致组织缺氧，严重时致死。这就是"煤气中毒"。
\end{thinkbox}

\begin{practice}
1. 空气中体积分数第二大的气体是\_\_\_\_\_。\\
2. 写出过氧化氢制氧气的化学方程式。\\
3. 燃烧的三条件是\_\_\_\_\_、\_\_\_\_\_、\_\_\_\_\_。\\
4. （判断）用带火星的木条伸入集气瓶——木条熄灭，说明是氧气。（\quad）\\
5. 实验室加热高锰酸钾制氧气的化学方程式为\_\_\_\_\_。\\
6. （判断）铁丝在氧气中燃烧生成Fe$_2$O$_3$。（\quad）\\
7. 燃烧的三个条件缺一不可，灭火的原理是破坏\_\_\_\_\_个条件即可。
\end{practice}

\newpage
\section{\textcolor{orange!80!black}{第二课\quad 水与氢气}}

\subsection*{水的组成——电解实验}

电解水（加少量硫酸或氢氧化钠增强导电性）：\\
负极产生\textbf{氢气}H$_2$，正极产生\textbf{氧气}O$_2$。体积比：\textbf{H$_2$ : O$_2$ = 2 : 1}。

\formula{2\text{H}_2\text{O} \xrightarrow{\text{通电}} 2\text{H}_2\uparrow + \text{O}_2\uparrow}

宏观上水由氢元素和氧元素组成；微观上每个水分子由两个氢原子和一个氧原子构成。

\subsection*{氢气——最轻的气体}

H$_2$是密度\textbf{最小}的气体（可充气球）。氢气具有\textbf{可燃性}（燃烧产物只有水——是最清洁的燃料）和\textbf{还原性}（可从金属氧化物中"夺走"氧）。

检验氢气纯度很重要——氢气和空气混合后点燃会发生\textbf{爆炸}（所以充氦气的气球比氢气安全得多，尽管氦气略贵）。

\subsection*{硬水与软水}

含较多可溶性钙镁化合物的水是\textbf{硬水}（肥皂不易起泡，水壶易结垢）；含较少或不含的是\textbf{软水}。煮沸可以降低水的硬度（暂时硬度）。

\subsection*{水的净化}

自来水厂的净化流程：\textbf{沉淀 $\rightarrow$ 过滤 $\rightarrow$ 吸附（活性炭）$\rightarrow$ 消毒（氯气或紫外线）}。蒸馏是最彻底的净化方法——获得纯水。

\begin{example}
电解水实验中，正极产生8 mL气体。负极产生多少mL气体？分别是什么？

\textbf{解：}正极：负极气体体积 = 1 : 2。正极8 mL $\rightarrow$ 负极16 mL。正极O$_2$，负极H$_2$。
\end{example}

\begin{practice}
1. 电解水时，正极产生\_\_\_\_\_，负极产生\_\_\_\_\_。正负极气体体积比\_\_\_\_\_。\\
2. （选择）下列关于硬水的说法正确的是（\quad）\\
\quad A. 硬水是纯水 \quad B. 煮沸可降低暂时硬度 \quad C. 硬水和软水用肉眼可区分\\
3. 氢气燃烧的化学方程式：\_\_\_\_\_。\\
4. （判断）电解水时正极产生氢气，负极产生氧气。（\quad）\\
5. 将硬水软化的常用家庭方法是\_\_\_\_\_。\\
6. 自来水净化流程中起吸附作用的物质是\_\_\_\_\_。\\
7. （计算）18g水完全电解，产生氢气的质量是多少？\_\_\_\_\_g。
\end{practice}

\newpage
\section{\textcolor{orange!80!black}{第三课\quad 碳和碳的氧化物}}

\subsection*{碳的单质——同一种元素，不同的"长相"}

同种元素可以形成不同的单质——这叫\textbf{同素异形体}。碳有几种著名的同素异形体：

\textbf{金刚石：}最硬的天然物质，无色透明晶体。每个碳原子与周围4个碳原子形成正四面体结构。用于切割玻璃、钻探。

\textbf{石墨：}深灰色、有金属光泽、滑腻。碳原子层状排列——层与层之间容易滑动（所以石墨能做铅笔芯和润滑剂）。石墨能导电——可用作电极。

\textbf{富勒烯（C$_{60}$等）和石墨烯：}现代纳米材料。C$_{60}$分子像足球（60个碳原子构成的球）。石墨烯是单层石墨——极薄、极强、导电性极好。

\subsection*{一氧化碳CO——有毒气体}

CO是碳不完全燃烧的产物。无色无味，有剧毒（与血红蛋白结合力极强）。CO具有\textbf{可燃性}（本身可燃烧生成CO$_2$）和\textbf{还原性}（可以还原金属氧化物，铁矿石炼铁就用CO）。

\subsection*{二氧化碳CO$_2$——温室气体}

\textbf{实验室制取CO$_2$：}CaCO$_3$ + 2HCl $\rightarrow$ CaCl$_2$ + H$_2$O + CO$_2$$\uparrow$（石灰石和大理石与稀盐酸反应）。用\textbf{向上排空气法}收集（CO$_2$密度大于空气）。

\textbf{CO$_2$的检验：}通入澄清石灰水——石灰水变\textbf{浑浊}：Ca(OH)$_2$ + CO$_2$ $\rightarrow$ CaCO$_3$$\downarrow$ + H$_2$O。这是检验CO$_2$的专属反应。

CO$_2$本身不支持燃烧也不可燃，可用于灭火。干冰（固态CO$_2$）升华时吸热——用于人工降雨和舞台烟雾。

\begin{example}
写出木炭还原氧化铜的化学方程式。

\textbf{解：}C + 2CuO $\xrightarrow{\text{高温}}$ 2Cu + CO$_2$$\uparrow$。\\
现象：黑色粉末变为\textbf{红色}（生成了铜），生成的气体使澄清石灰水变浑浊。
\end{example}

\begin{practice}
1. 金刚石和石墨物理性质差异大的原因是\_\_\_\_\_不同。\\
2. 检验CO$_2$的方法是\_\_\_\_\_。\\
3. 写出CaCO$_3$与稀HCl反应的化学方程式。\\
4. （判断）CO和CO$_2$都有毒。（\quad）\\
5. 写出CO燃烧的化学方程式\_\_\_\_\_。\\
6. （选择）下列物质中碳元素化合价最高的是（\quad）\\
\quad A. CO \quad B. CO$_2$ \quad C. CH$_4$ \quad D. C\\
7. 实验室制取CO$_2$不用稀硫酸是因为\_\_\_\_\_。
\end{practice}

\newpage
\section{\textcolor{orange!80!black}{第四课\quad 金属}}

\subsection*{金属的共性}

大多数金属具有：金属光泽、导电导热性、延展性（可以拉成丝、压成薄片）。汞（水银）是唯一的液态金属。

\subsection*{金属活动性顺序——"谁更活泼"}

\begin{center}
\textbf{K Ca Na Mg Al Zn Fe Sn Pb (H) Cu Hg Ag Pt Au}\\
\small{（从强到弱：钾钙钠镁铝锌铁锡铅\ 氢\ 铜汞银铂金）}
\end{center}

\textbf{H之前}的金属能从酸中置换出氢气；\textbf{排在前面的金属}能把排在后面的金属从其盐溶液中置换出来（"强置换弱"）。

\subsection*{铁的冶炼——高炉炼铁}

赤铁矿（Fe$_2$O$_3$）或磁铁矿（Fe$_3$O$_4$）在高炉中用CO还原：

\formula{\text{Fe}_2\text{O}_3 + 3\text{CO} \xrightarrow{\text{高温}} 2\text{Fe} + 3\text{CO}_2}

\subsection*{铝的"自我保护"——氧化膜}

铝的化学性质其实很活泼——但铝制品在空气中会迅速形成一层致密的\textbf{氧化铝Al$_2$O$_3$薄膜}。这层膜像一层透明的"铠甲"，阻止了内部的铝继续氧化。这就是为什么铝制品不容易生锈——其实它已经"锈"了一薄层，但那层"锈"恰好保护了它。

\textbf{铁生锈}则相反——生成的铁锈（Fe$_2$O$_3$·xH$_2$O）疏松多孔，不能阻止内部继续氧化。铁锈"海绵吸水"般的结构让氧气和水不断渗入——加速腐蚀。

\begin{example}
判断下列反应能否发生：① Fe + CuSO$_4$ $\rightarrow$ ？ ② Cu + FeSO$_4$ $\rightarrow$ ？

\textbf{解：}① Fe在Cu之前 $\rightarrow$ 能置换Cu：Fe + CuSO$_4$ $\rightarrow$ FeSO$_4$ + Cu。\\
② Cu在Fe之后 $\rightarrow$ 不能置换Fe，无反应。
\end{example}

\begin{practice}
1. 金属活动性顺序中，\_\_\_\_\_是最活泼的金属（常见）。\\
2. 写出铝和稀盐酸反应的化学方程式。\\
3. （判断）Cu可以从AgNO$_3$溶液中置换出Ag。（\quad）\\
4. 铁生锈的条件是铁与\_\_\_\_\_和\_\_\_\_\_同时接触。\\
5. （判断）铝制品耐腐蚀是因为铝的化学性质不活泼。（\quad）\\
6. 写出铁和硫酸铜溶液反应的化学方程式\_\_\_\_\_。\\
7. 写出一种防止铁生锈的措施\_\_\_\_\_。
\end{practice}

\newpage
\section{\textcolor{orange!80!black}{第五课\quad 溶液与溶解度}}

\subsection*{溶液——均匀、稳定、混合物}

溶液由\textbf{溶质}（被溶解的物质）和\textbf{溶剂}（溶解溶质的物质，通常是水）组成。溶液是均一、稳定的混合物。

\subsection*{饱和与不饱和}

一定温度下，在一定量溶剂中，\textbf{不能再溶解某种溶质的溶液}叫该溶质的\textbf{饱和溶液}。还能继续溶解的叫\textbf{不饱和溶液}。

\subsection*{溶解度}

\begin{definition}{溶解度}
在一定温度下，某固体物质在100g溶剂中达到饱和状态时所溶解的质量，叫做该物质在这种溶剂中的\textbf{溶解度}。单位：g。
\end{definition}

大多数固体物质的溶解度随温度升高而\textbf{增大}（如KNO$_3$——溶解度曲线陡然上升），少数变化不大（如NaCl），极少数反而降低（如Ca(OH)$_2$）。

\subsection*{结晶}

\textbf{降温结晶（冷却热饱和溶液）：}溶解度随温度显著变化的物质适用（如硝酸钾KNO$_3$）。\\
\textbf{蒸发结晶：}溶解度随温度变化不大的物质适用（如食盐NaCl）。

\begin{example}
在20℃时，36g NaCl恰好溶于100g水形成饱和溶液。求20℃时NaCl的溶解度。

\textbf{解：}溶解度 = 36g（每100g水最多溶解36g NaCl）。
\end{example}

\begin{practice}
1. 溶液由\_\_\_\_\_和\_\_\_\_\_组成。\\
2. 大多数固体物质的溶解度随温度升高而\_\_\_\_\_。\\
3. （判断）饱和溶液一定是浓溶液。（\quad）\\
4. 从KNO$_3$溶液中获得KNO$_3$晶体，宜用\_\_\_\_\_结晶法。\\
5. 20℃时NaCl的溶解度为36g，该温度下50g水中最多溶解\_\_\_\_\_g NaCl。\\
6. （判断）搅拌可以增大固体物质在水中的溶解度。（\quad）\\
7. 从海水中提取食盐常用\_\_\_\_\_结晶法。
\end{practice}

\newpage
\section{\textcolor{orange!80!black}{第六课\quad 常见的酸和碱}}

\subsection*{酸——氢离子H$^+$的世界}

\begin{definition}{酸}
电离时生成的阳离子\textbf{全部是H$^+$}的化合物。
\end{definition}

\textbf{常见酸：}盐酸HCl、硫酸H$_2$SO$_4$、硝酸HNO$_3$、碳酸H$_2$CO$_3$（易分解）。

\textbf{酸的共性（H$^+$的性质）：}
\begin{itemize}
  \item 能使紫色石蕊试液变\textbf{红}
  \item 能和活泼金属（H之前）反应产生H$_2$
  \item 能和金属氧化物反应生成盐和水
  \item 能和碱发生\textbf{中和反应}生成盐和水
\end{itemize}

\subsection*{碱——氢氧根OH$^-$的世界}

\begin{definition}{碱}
电离时生成的阴离子\textbf{全部是OH$^-$}的化合物。
\end{definition}

\textbf{常见碱：}氢氧化钠NaOH（烧碱/火碱，强腐蚀性）、氢氧化钙Ca(OH)$_2$（熟石灰/消石灰）。

\textbf{碱的共性（OH$^-$的性质）：}
\begin{itemize}
  \item 能使紫色石蕊试液变\textbf{蓝}，使无色酚酞变\textbf{红}
  \item 能和非金属氧化物（如CO$_2$）反应生成盐和水
  \item 能和酸发生中和反应
\end{itemize}

\subsection*{pH——酸碱性的一把"尺子"}

\formula{\text{pH} < 7\text{：酸性} \qquad \text{pH} = 7\text{：中性} \qquad \text{pH} > 7\text{：碱性}}

pH越小酸性越强；pH越大碱性越强。用pH试纸可以粗略测定溶液的pH。

\subsection*{中和反应——酸和碱的"和解"}

酸 + 碱 $\rightarrow$ 盐 + 水。例：HCl + NaOH $\rightarrow$ NaCl + H$_2$O。中和反应在农业（改良酸性土壤用熟石灰）、工业（处理酸性废水）和医疗（用含Al(OH)$_3$的胃药中和过多胃酸）中都有应用。

\begin{example}
写出稀硫酸和氢氧化钠反应的化学方程式。

\textbf{解：}H$_2$SO$_4$ + 2NaOH $\rightarrow$ Na$_2$SO$_4$ + 2H$_2$O（中和反应）。
\end{example}

\begin{practice}
1. 酸的水溶液中共同的阳离子是\_\_\_\_\_；碱的共同阴离子是\_\_\_\_\_。\\
2. pH=3的溶液呈\_\_\_\_\_性，pH=11的溶液呈\_\_\_\_\_性。\\
3. （选择）下列物质中属于碱的是（\quad）A. NaCl B. NaOH C. HCl D. H$_2$O\\
4. 写出盐酸和氢氧化钙的中和反应方程式。\\
5. （选择）下列物质中pH最小的是（\quad）\\
\quad A. 食醋 \quad B. 纯水 \quad C. 肥皂水 \quad D. 石灰水\\
6. 写出稀硫酸和氧化铁反应的化学方程式\_\_\_\_\_。\\
7. （判断）所有金属都能与稀盐酸反应生成氢气。（\quad）
\end{practice}

\newpage
\section{\textcolor{orange!80!black}{第七课\quad 盐与复分解反应}}

\subsection*{盐——金属离子（或NH$_4^+$）+ 酸根离子}

酸和碱中和的产物就是盐。NaCl（食盐）、CaCO$_3$（石灰石的主要成分）、Na$_2$CO$_3$（纯碱/苏打）、NaHCO$_3$（小苏打）都是盐。

\subsection*{复分解反应}

\begin{definition}{复分解反应}
两种化合物互相交换成分，生成另外两种化合物的反应。\\
通式：\textbf{AB + CD $\rightarrow$ AD + CB}。
\end{definition}

\textbf{复分解反应发生的条件（生成物至少满足一条）：}生成\textbf{沉淀}、或生成\textbf{气体}、或生成\textbf{水}。

\subsection*{盐的溶解性口诀}

钾钠铵盐全都溶（K$^+$ Na$^+$ NH$_4^+$的盐都可溶），硝酸盐遇水影无踪（NO$_3^-$盐都可溶）。\\
盐酸不溶银亚汞（AgCl、Hg$_2$Cl$_2$不溶），硫酸不溶钡和铅（BaSO$_4$、PbSO$_4$不溶）。\\
碳酸只溶钾钠铵（大多数碳酸盐不溶）。

\begin{example}
判断下列反应能否发生：NaCl + AgNO$_3$ $\rightarrow$ ？

\textbf{解：}NaCl + AgNO$_3$ $\rightarrow$ NaNO$_3$ + AgCl$\downarrow$。AgCl为不溶沉淀——反应\textbf{能发生}。
\end{example}

\begin{practice}
1. 复分解反应发生的条件是生成\_\_\_\_\_、\_\_\_\_\_或\_\_\_\_\_。\\
2. （判断）NaCl和KNO$_3$混合能发生复分解反应。（\quad）\\
3. 写出Na$_2$CO$_3$与稀HCl反应的化学方程式。\\
4. 写出一种不溶于水的硫酸盐：\_\_\_\_\_。\\
5. （判断）AgCl不溶于水，也不溶于稀硝酸。（\quad）\\
6. 写出CaCO$_3$与稀HCl反应的离子方程式\_\_\_\_\_。\\
7. 写出一种常见的氮肥\_\_\_\_\_。
\end{practice}

\newpage
\section{\textcolor{orange!80!black}{第八课\quad 氧化还原反应（高中深化）}}

\subsection*{从"得氧失氧"到"化合价升降"}

初中阶段，氧化反应 = 得氧，还原反应 = 失氧。这个定义太窄了。高中阶段用\textbf{化合价变化}来定义：

\begin{definition}{氧化还原反应}
\textbf{氧化反应}：物质所含元素化合价\textbf{升高}（失去电子）的过程。\\
\textbf{还原反应}：物质所含元素化合价\textbf{降低}（得到电子）的过程。\\
氧化反应和还原反应\textbf{同时发生}——有升必有降。
\end{definition}

\textbf{氧化剂}：在反应中\textbf{化合价降低}（得电子）的物质——被还原。\\
\textbf{还原剂}：在反应中\textbf{化合价升高}（失电子）的物质——被氧化。

\subsection*{化合价升降法配平}

这是配平氧化还原反应最系统的方法：
\begin{enumerate}
  \item 标出变价元素的化合价
  \item 算出每原子化合价的变化值
  \item 乘以原子个数得总变化值
  \item 求最小公倍数，确定氧化剂和还原剂的系数
  \item 配平其他元素（H、O最后调）
\end{enumerate}

\begin{example}
用化合价升降法配平：Cu + HNO$_3$（稀）$\rightarrow$ Cu(NO$_3$)$_2$ + NO$\uparrow$ + H$_2$O

\textbf{解：}Cu: 0 $\rightarrow$ +2（升2）×3；N: +5 $\rightarrow$ +2（降3）×2。\\
最小公倍数6——Cu系数3，HNO$_3$作氧化剂的部分系数2（还有作酸的HNO$_3$）。\\
\textbf{3Cu + 8HNO$_3$（稀）$\rightarrow$ 3Cu(NO$_3$)$_2$ + 2NO$\uparrow$ + 4H$_2$O}。
\end{example}

\begin{thinkbox}
高锰酸钾KMnO$_4$是极强的氧化剂（Mn为+7价），在反应中常常被还原成Mn$^{2+}$（+2价，溶液从紫色变为无色/浅粉）。生活中用高锰酸钾溶液泡脚（除脚气）就是利用了它的氧化性——氧化杀灭细菌。
\end{thinkbox}

\begin{practice}
1. （判断）氧化还原反应的本质是氧的得失。（\quad）\\
2. 指出反应 Fe + CuSO$_4$ $\rightarrow$ FeSO$_4$ + Cu 中的氧化剂和还原剂。\\
3. （填空）氧化剂在反应中化合价\_\_\_\_\_（升高/降低），被\_\_\_\_\_（氧化/还原）。\\
4. 配平：C + HNO$_3$ $\rightarrow$ CO$_2$ + NO$_2$ + H$_2$O\\
5. 在反应2H$_2$ + O$_2$ $\rightarrow$ 2H$_2$O中，氧化剂是\_\_\_\_\_，还原剂是\_\_\_\_\_。\\
6. （判断）有单质生成的反应一定是氧化还原反应。（\quad）\\
7. 配平：Fe$_2$O$_3$ + CO $\rightarrow$ Fe + CO$_2$
\end{practice}

\newpage
\section{\textcolor{orange!80!black}{第九课\quad 离子反应}}

\subsection*{电解质与非电解质}

\begin{definition}{电解质}
在水溶液中或熔融状态下能导电的化合物（酸、碱、盐）。导电是因为存在\textbf{自由移动的离子}。
\end{definition}

\textbf{强电解质：}完全电离——强酸（HCl、HNO$_3$、H$_2$SO$_4$）、强碱（NaOH、KOH）、大多数盐。\\
\textbf{弱电解质：}部分电离——弱酸（CH$_3$COOH）、弱碱（NH$_3$·H$_2$O）、水。

\subsection*{离子反应与离子方程式}

有些反应的本质不是"分子之间"的反应，而是\textbf{离子之间}的反应。用离子方程式来表示：

\begin{enumerate}
  \item 把强电解质拆成离子形式
  \item 弱电解质、沉淀、气体、水——保留化学式
  \item 删除两边都不参与反应相同的离子（"旁观离子"）
  \item 检查电荷守恒和原子守恒
\end{enumerate}

\begin{example}
写出盐酸和氢氧化钠反应的离子方程式。

\textbf{解：}化学方程式：HCl + NaOH $\rightarrow$ NaCl + H$_2$O。\\
拆离子：H$^+$ + Cl$^-$ + Na$^+$ + OH$^-$ $\rightarrow$ Na$^+$ + Cl$^-$ + H$_2$O。\\
消去Na$^+$和Cl$^-$（旁观离子）：\textbf{H$^+$ + OH$^-$ $\rightarrow$ H$_2$O}。\\
这才是中和反应的本质——酸中的H$^+$和碱中的OH$^-$结合成水。
\end{example}

\begin{example}
写出Na$_2$CO$_3$溶液与稀盐酸反应的离子方程式。

\textbf{解：}CO$_3^{2-}$ + 2H$^+$ $\rightarrow$ H$_2$O + CO$_2$$\uparrow$（CO$_3^{2-}$与H$^+$反应生成碳酸，碳酸立即分解成水和CO$_2$）。
\end{example}

\begin{thinkbox}
为什么同样的"酸+碱"反应（HCl+NaOH、HNO$_3$+KOH等）写出的离子方程式都一样——H$^+$+OH$^-$$\rightarrow$H$_2$O？因为中和反应的本质不关心酸和碱的"名字"——只关心H$^+$和OH$^-$的结合。化学方程式告诉了你"具体是谁和谁"，离子方程式告诉你的是"本质上发生了什么"。
\end{thinkbox}

\begin{practice}
1. （判断）电解质在水溶液中全部以离子形式存在。（\quad）\\
2. 写出AgNO$_3$溶液与NaCl溶液反应的离子方程式。\\
3. 下列不能拆成离子的是（\quad）\\
\quad A. HCl \quad B. NaCl \quad C. CaCO$_3$（固体）\quad D. NaOH\\
4. 离子方程式两侧必须同时满足\_\_\_\_\_守恒和\_\_\_\_\_守恒。\\
5. （判断）BaSO$_4$难溶于水，所以它是弱电解质。（\quad）\\
6. 写出NaOH与稀H$_2$SO$_4$反应的离子方程式\_\_\_\_\_。\\
7. 在书写离子方程式时，\_\_\_\_\_、\_\_\_\_\_、\_\_\_\_\_等物质不能拆成离子。
\end{practice}

\vspace{0.5cm}
\begin{center}\rule{0.7\textwidth}{1pt}\vspace{0.4cm}
{\large\textcolor{blue!70!black}{\textbf{—— 化学2·常见元素与无机反应 完 ——}}}
\vspace{0.3cm}\textcolor{gray}{\small 下一本预告：化学3·化学反应原理——平衡、速率与电化学}
\end{center}

\newpage

\appendix
\section*{\textcolor{blue!70!black}{附录A\quad 元素周期表}}

\begin{figplaceholder}
此处插入完整元素周期表（118元素，彩色，含中文名称与相对原子质量）。\\
本图化学1-4通用（见 figure/requirements/periodic-table.txt）
\end{figplaceholder}

\subsection*{1-20号元素速查表}

\begin{center}
\begin{tabular}{|c|c|c|c|c|}
\hline
\textbf{序数} & \textbf{符号} & \textbf{名称} & \textbf{相对原子质量} & \textbf{常见化合价} \\
\hline
1 & H & 氢 & 1.008 & +1 \\
2 & He & 氦 & 4.003 & 0 \\
3 & Li & 锂 & 6.941 & +1 \\
4 & Be & 铍 & 9.012 & +2 \\
5 & B & 硼 & 10.81 & +3 \\
6 & C & 碳 & 12.01 & +2,+4,-4 \\
7 & N & 氮 & 14.01 & -3,+3,+5 \\
8 & O & 氧 & 16.00 & -2 \\
9 & F & 氟 & 19.00 & -1 \\
10 & Ne & 氖 & 20.18 & 0 \\
11 & Na & 钠 & 22.99 & +1 \\
12 & Mg & 镁 & 24.31 & +2 \\
13 & Al & 铝 & 26.98 & +3 \\
14 & Si & 硅 & 28.09 & +4,-4 \\
15 & P & 磷 & 30.97 & -3,+3,+5 \\
16 & S & 硫 & 32.07 & -2,+4,+6 \\
17 & Cl & 氯 & 35.45 & -1,+1,+5,+7 \\
18 & Ar & 氩 & 39.95 & 0 \\
19 & K & 钾 & 39.10 & +1 \\
20 & Ca & 钙 & 40.08 & +2 \\
\hline
\end{tabular}
\end{center}

\vspace{0.3cm}
\textcolor{gray}{\small 注：相对原子质量为近似值，教学用。稀有气体（He/Ne/Ar）化合价为0。}

\newpage
\section*{\textcolor{defblue}{参考答案}}

\subsection*{第一课}
1. 氧气（O$_2$）\\
2. 2H$_2$O$_2$ $\xrightarrow{\text{MnO}_2}$ 2H$_2$O + O$_2$$\uparrow$\\
3. 可燃物、氧气（空气）、达到着火点\\
4. （×）\\
5. 2KMnO$_4 \xrightarrow{\Delta}$ K$_2$MnO$_4$ + MnO$_2$ + O$_2$$\uparrow$\\
6. （×）生成Fe$_3$O$_4$\\
7. 任意一个

\subsection*{第二课}
1. 氧气（O$_2$）、氢气（H$_2$），1:2\\
2. B\\
3. 2H$_2$ + O$_2$ $\xrightarrow{\text{点燃}}$ 2H$_2$O\\
4. （×）正极产生氧气，负极产生氢气\\
5. 煮沸\\
6. 活性炭\\
7. 2g（$18\text{g} \times \frac{2}{18} = 2\text{g}$）

\subsection*{第三课}
1. 碳原子的排列方式\\
2. 通入澄清石灰水——石灰水变浑浊\\
3. CaCO$_3$ + 2HCl $\rightarrow$ CaCl$_2$ + H$_2$O + CO$_2$$\uparrow$\\
4. （×）CO$_2$无毒，CO有毒\\
5. 2CO + O$_2$ $\xrightarrow{\text{点燃}}$ 2CO$_2$\\
6. B（CO$_2$中C为+4价）\\
7. 生成的CaSO$_4$微溶，覆盖在大理石表面阻止反应继续

\subsection*{第四课}
1. 钾（K）\\
2. 2Al + 6HCl $\rightarrow$ 2AlCl$_3$ + 3H$_2$$\uparrow$\\
3. （√）\\
4. 氧气（O$_2$）和水（H$_2$O）\\
5. （×）铝表面形成致密氧化膜阻止继续氧化\\
6. Fe + CuSO$_4$ $\rightarrow$ FeSO$_4$ + Cu\\
7. 涂油漆、镀锌、制成不锈钢等（合理即可）

\subsection*{第五课}
1. 溶质、溶剂\\
2. 增大\\
3. （×）\\
4. 降温结晶（冷却热饱和溶液）\\
5. 18（$36 \times 50/100 = 18$）\\
6. （×）搅拌只加快溶解速率，不改变溶解度\\
7. 蒸发

\subsection*{第六课}
1. H$^+$、OH$^-$\\
2. 酸、碱\\
3. B\\
4. 2HCl + Ca(OH)$_2$ $\rightarrow$ CaCl$_2$ + 2H$_2$O\\
5. A（食醋含醋酸呈酸性，pH最小）\\
6. H$_2$SO$_4$ + Fe$_2$O$_3$ $\rightarrow$ Fe$_2$(SO$_4$)$_3$ + 3H$_2$O\\
7. （×）铜、银等不活泼金属不与稀盐酸反应

\subsection*{第七课}
1. 沉淀、气体、水\\
2. （×）\\
3. Na$_2$CO$_3$ + 2HCl $\rightarrow$ 2NaCl + H$_2$O + CO$_2$$\uparrow$\\
4. BaSO$_4$（或PbSO$_4$）\\
5. （√）\\
6. CaCO$_3$ + 2H$^+$ $\rightarrow$ Ca$^{2+}$ + H$_2$O + CO$_2$$\uparrow$\\
7. 尿素（NH$_2$)$_2$CO、硝酸铵NH$_4$NO$_3$等（合理即可）

\subsection*{第八课}
1. （×）本质是电子的得失（化合价变化）\\
2. 氧化剂：CuSO$_4$（Cu$^{2+}$），还原剂：Fe\\
3. 降低、还原\\
4. C + 4HNO$_3$ $\rightarrow$ CO$_2$ + 4NO$_2$ + 2H$_2$O\\
5. 氧化剂：O$_2$，还原剂：H$_2$\\
6. （×）同素异形体转化（如$3\text{O}_2 \rightarrow 2\text{O}_3$）不是氧化还原\\
7. Fe$_2$O$_3$ + 3CO $\rightarrow$ 2Fe + 3CO$_2$

\subsection*{第九课}
1. （×）弱电解质在水中部分电离\\
2. Ag$^+$ + Cl$^-$ $\rightarrow$ AgCl$\downarrow$\\
3. C\\
4. 原子（质量）守恒、电荷守恒\\
5. （×）BaSO$_4$难溶但溶于水的部分完全电离，属于强电解质\\
6. H$^+$ + OH$^-$ $\rightarrow$ H$_2$O\\
7. 沉淀、气体、弱电解质（或水）

\end{document}
